\documentclass{article}

\textwidth=6.5in
\textheight=8.5in
\voffset=-54pt
\oddsidemargin=0in
\evensidemargin=0in
\setlength{\parskip}{8pt}
\setlength{\parindent}{0pt}

\usepackage{workbook}

\usepackage[sols]{handout}

\begin{document}

\begin{center}
  \large\textsc{Problem Set 5 - Logarithmic Differentiation}
  
      \pspace{12pt}
  \begin{problemonly}\large Name: \rule{5cm}{0.15mm}\\ \end{problemonly}
\end{center}

\begin{grading}
\begin{enumerate}
    \item 12 points
    \item 10 points
    \item 5 points
    \item 20 points
    \item 3 points
\end{enumerate}
Total: 50 points
\end{grading}


% \begin{problemonly}
%   Technically, you should only use logarithmic differentiation on positive
%   functions, since you can only take the log of a positive quantity.  The
%   functions on this homework are all positive, so you can use logarithmic
%   differentation without fear!
% \end{problemonly}

\begin{enumerate}[topsep=0pt]
\item

\begin{grading}
12 points. 4 points for each part.
\end{grading}

 You may find examples 17.2 and
  17.3 in Gottlieb helpful for the following
  problems.

 You need not simplify your
  answers.
  \begin{enumerate}

  \item % 17.2.1(b)
    \begin{problem}[\vfill]
      If $f(x) =5(x^2+1)^x$, find $f'(x)$.
    \end{problem}

    \begin{solution}
      Since the base and exponent of $(x^2 + 1)^x$ both depend on $x$, we
      need logarithmic differentiation.  You can do logarithmic
      differentiation on just $(x^2 + 1)^x$, or you can do it on $f(x) = 5
      (x^2 + 1)^x$; here, we'll do the latter:
      \begin{align*}
        f(x) &= 5(x^2 + 1)^x \\
        \intertext{Take the natural log of both sides and simplify the
          right side:}
        \ln f(x) &= \ln [5 (x^2 + 1)^x] \\
        &= \ln 5 + \ln \left[ (x^2 + 1)^x \right] \\
        &= \ln 5 + x \ln (x^2 + 1) \\
        \intertext{Differentiate both sides with respect to $x$:}
        \frac{1}{f(x)} f'(x) &= 0 + \frac{d}{dx}
        \left(\fcolorbox{green}{white}{$\displaystyle x$}\cdot
          \fcolorbox{green}{white}{$\displaystyle \ln\left(x^
              { 2} + 1\right)$}\right) \\
        &= x\cdot \frac{d}{dx} \left(\fcolorbox{red}{white}{$\displaystyle
            \ln\left(\fcolorbox{blue}{white}{$\displaystyle x^ { 2} +
                1$}\right)$}\right)+ \ln\left(x^ { 2} + 1\right) \\
        &= x\cdot \frac{1}{x^ { 2} + 1}\cdot 2
        x+ \ln\left(x^ { 2} + 1\right) \\
        &= \frac{2x^2}{x^2 + 1} + \ln\left(x^ { 2} + 1\right) \\
        \shortintertext{Solve for $f'(x)$:} f'(x) &= f(x) \left(
          \frac{2x^2}{x^2 + 1} + \ln\left(x^ { 2} + 1\right) \right) \\
        &= \boxed{5(x^2 + 1)^x \left( \frac{2x^2}{x^2 + 1} + \ln\left(x^ {
                2} + 1\right) \right)}
      \end{align*}
    \end{solution}

  \item % 17.2.3(b)
    \begin{problem}[\vfill\newpage]
      If $y = \dfrac{2^t}{t^{2t}}$, find
      $\dfrac{dy}{dt}$.
    \end{problem}

    \begin{solution}
      We need logarithmic differentiation to differentiate $t^{2t}$, so we
      might as well use it on $y$ (since we'll be able to use log rules to
      simplify $\ln y$.  As always, we start with the given equation:
      \begin{align*}
        y &= \frac{2^t}{t^{2t}} \\
        \intertext{Take the natural log of both sides and simplify the
          right side:}
        \ln y &= \ln \left( \frac{2^t}{t^{2t}} \right) \\
        &= \ln (2^t) - \ln (t^{2t}) \\
        &= t \ln 2 - 2t \ln t \\
        \intertext{Differentiate both sides with respect to $t$:}
        \frac{1}{y} \frac{dy}{dt} &=
        \ln 2-2\cdot \frac{d}{dt} \left(\fcolorbox{green}{white}{$\displaystyle  t$}\cdot \fcolorbox{green}{white}{$\displaystyle  \ln t$}\right) \\
        &= \ln 2 -2 \left(t\cdot \frac{1}{t}+ \ln t\right) \\
        &= \ln 2 - 2 - 2 \ln t \\
        \shortintertext{Solve for $\frac{dy}{dt}$:}
        \frac{dy}{dt} &= \boxed{\frac{2^t}{t^{2t}} \left( \ln 2 - 2 - 2
            \ln t \right)}
      \end{align*}    
    \end{solution}

  \item % 17.2.4(b), modified
    \begin{problem}[\vfill]
      If $y = \displaystyle\frac{xe^{5x}}{(3x+1)^2 \sqrt{x-2}}$, find
      $\dfrac{dy}{dx}$ using logarithmic differentiation.
    \end{problem}

    \begin{solution}
      We start with the given equation:
      \begin{align*}
        y &= \frac{xe^{5x}}{(3x+1)^2 \sqrt{x-2}} \\
        \intertext{Take the natural log of both sides and simplify the
          right side:}
        \ln y &= \ln \left( \frac{xe^{5x}}{(3x+1)^2 \sqrt{x-2}} \right) \\
        &= \ln \left( x e^{5x} \right) - \ln \left[ (3x+1)^2 \sqrt{x-2}
        \right] \\
        &= \ln x + \ln \left( e^{5x} \right) - \ln \left[ (3x + 1)^2
        \right] - \ln \sqrt{x - 2} \\
        &= \ln x + 5x - 2 \ln (3x + 1) - \frac{1}{2} \ln (x - 2)  \\
        \intertext{Differentiate both sides with respect to $x$:}
        \frac{1}{y} \frac{dy}{dx} &= \frac{1}{x} + 5 - \frac{2}{3x + 1}
        \cdot 3 - \frac{1}{2} \cdot \frac{1}{x - 2} \cdot 1 \\
        &= \frac{1}{x} + 5 - \frac{6}{3x + 1} - \frac{1}{2 (x - 2)} \\
        \shortintertext{Solve for $\frac{dy}{dx}$:} \frac{dy}{dx} &=
        \boxed{\frac{xe^{5x}}{(3x+1)^2 \sqrt{x-2}} \left( \frac{1}{x} + 5
            - \frac{6}{3x + 1} - \frac{1}{2 (x - 2)} \right)}
      \end{align*}
    \end{solution}
  \end{enumerate}





\item

\begin{grading}
10 points. 5 points for each part.
\end{grading}

Find $f'(x)$. You need not simplify your answers.  You will
  need to use logarithmic differentiation on a part of the formula for $f(x)$, not the entire expression.
  You may find it useful to give that part a name, like $g(x)$ or $z$,
  then work on finding $g'(x)$ or $\frac{dz}{dx}$ separately. (Why does it
  not work to use logarithmic differentiation on the entire formula?)
  \begin{enumerate}
  \item
    \begin{problem}[\vfill\newpage]
      $f(x) =x(2x^3 +1)^x +5$
    \end{problem}

    \begin{solution}
      It won't work to do logarithmic differentiation directly on $f$
      because we can't simplify the natural log of a sum.  But the piece
      of $f$ that we really need logarithmic differentiation for is the
      first one, so let's give that a name.  If we say $g(x) = x(2x^3
      +1)^x$, then $f(x) = g(x) + 5$, so $f'(x) = g'(x) + 0$.  We'll use
      logarithmic differentiation to find $g'(x)$:
      \begin{align*}
        g(x) &= x (2x^3 + 1)^x \\
        \intertext{Take the natural log of both sides and simplify the
          right side:}
        \ln g(x) &= \ln \left[ x (2x^3 + 1)^x \right] \\
        &= \ln x + \ln \left[ (2x^3 + 1)^x \right] \\
        &= \ln x + x \ln \left( 2x^3 + 1 \right) \\
        \intertext{Differentiate both sides with respect to $x$:}
        \frac{1}{g(x)} g'(x) &= \frac{1}{x}+ \left(x\cdot \frac{d}{dx}
          \left(\fcolorbox{red}{white}{$\displaystyle
              \ln\left(\fcolorbox{blue}{white}{$\displaystyle 2 x^ { 3} +
                  1$}\right)$}\right)+ \ln\left(2 x^ { 3} +
            1\right)\right) \\
        &= \frac{1}{x}+ \left(x\cdot \frac{1}{2 x^ { 3} + 1}\cdot 6 x^ {
            2} + \ln\left(2 x^ { 3} +
            1\right)\right) \\
        &= \frac{1}{x}+ \frac{6x^3}{2 x^ { 3} + 1} + \ln\left(2 x^ { 3} +
          1\right) \\
        \shortintertext{Solve for $g'(x)$:}
        g'(x) &= x (2x^3 + 1)^x \left( \frac{1}{x}+ \frac{6x^3}{2 x^ { 3}
            + 1} + \ln\left(2 x^ { 3} + 
            1\right) \right) 
      \end{align*}
      Finally, remember that we already said that $f'(x) = g'(x)$, so
      \[f'(x) = \boxed{x (2x^3 + 1)^x \left( \frac{1}{x}+ \frac{6x^3}{2 x^
            { 3} + 1} + \ln\left(2 x^ { 3} + 1\right) \right)}.\]
    \end{solution}

  \item
\begin{problem}[\vfill]
      $f(x) = 3\cdot 2^x +2\cdot x^3 +3\cdot x^{2x+3}$
    \end{problem}

    \begin{solution}
      The term we need logarithmic differentiation for is $3 \cdot x^{2x +
        3}$ (or just $x^{2x + 3}$), so let $g(x) = x^{2x + 3}$.  Then,
      $f(x) = 3 \cdot 2^x + 2 \cdot x^3 + 3 \cdot g(x)$, so
      \begin{equation}
f'(x) = 3 \cdot 2^x \ln 2 + 6x^2 + 3 g'(x)
      \end{equation}
      We'll use logarithmic differentiation to find $g'(x)$:
      \begin{align*}
        g(x) &= x^{2x+3} \\
        \intertext{Take the natural log of both sides and simplify the
          right side:}
        \ln g(x) &= \ln \left( x^{2x + 3} \right) \\
        &= (2x + 3) \ln x \\
        \intertext{Differentiate both sides with respect to $x$:}
        \frac{1}{g(x)} g'(x) &= \frac{d}{dx}
        \left(\fcolorbox{green}{white}{$\displaystyle \left(2 x+
              3\right)$}\cdot
          \fcolorbox{green}{white}{$\displaystyle  \ln x$}\right) \\
        &= (2 x+ 3)\cdot \frac{1}{x}+ 2 \ln x \\
        \shortintertext{Solve for $g'(x)$:} g'(x) &= x^{2x + 3} \left[ (2
          x+ 3)\cdot \frac{1}{x}+ 2 \ln x \right]
      \end{align*}
      Then, by (1), $f'(x) = \boxed{3 \cdot 2^x
        \ln 2 + 6x^2 + 3 x^{2x + 3} \left[ (2 x+ 3)\cdot \frac{1}{x}+ 2
          \ln x \right]}$.      
    \end{solution}
  \end{enumerate}




\item 

\begin{grading}
5 points. Students will need to set the function equal to some variable in order to take the log of both sides of an equation.
\end{grading}

% 17.2.7
  \begin{problem}[\vfill]
    Come up with a differentiation rule (like the Product Rule) for $\diff*{\!\left[f(x)^{g(x)}\right]}{x}$. Your formula should give the derivative in terms of $f(x)$, $f'(x)$, $g(x)$, and $g'(x)$.
  \end{problem}

  \begin{solution}
    As usual, we should start by giving a name to the thing we're trying
    to differentiate using logarithmic differentiation; let's call it
    $h(x)$: 
    \begin{align*}
      h(x) &= f(x)^{g(x)} \\
      \intertext{Take $\ln$ of both sides:}
      \ln h(x) &= \ln \left( f(x)^{g(x)} \right) \\
      &= g(x) \ln f(x) \\
      \intertext{Differentiate both sides:}
      \frac{h'(x)}{h(x)} &= g(x) \frac{f'(x)}{f(x)} + g'(x) \ln f(x) \\
      \intertext{Solve for $h'(x)$:}
      h'(x) &= h(x) \left[ g(x) \frac{f'(x)}{f(x)} + g'(x) \ln f(x)
      \right] \\
      &= f(x)^{g(x)} \left[ g(x) \frac{f'(x)}{f(x)} + g'(x) \ln f(x)
      \right] 
    \end{align*}
  \end{solution}


\pspace{\newpage}

\item 

\begin{grading}
 20 points. Points are distributed to parts (a)-(g) as follows:\\
 4/4/2/2/4/2/2
\end{grading}

In this problem, you'll look at the function $f(x) = x^{\ln x}$ on
  its natural domain $(0, \infty)$.
  \begin{enumerate}
  \item
    \begin{problem}[\stretch{1}]
      Find $f'(x)$.
    \end{problem}

    \begin{solution}
      \ideas{Since the base and exponent both depend on $x$, we should use
        logarithmic differentiation.}  We start with the given equation:
      \begin{align*}
        f(x) &= x^{\ln x} \\
        \intertext{Take the natural log of both sides and simplify the
          right side:}
        \ln f(x) &= \ln \left( x^{\ln x} \right) \\
        &= (\ln x)(\ln x) \\
        &= (\ln x)^2 \\
        \intertext{Differentiate both sides with respect to $x$:}
        \frac{1}{f(x)} f'(x) &= 2 (\ln x) \frac{1}{x} \\
        \shortintertext{Solve for $f'(x)$:}
        f'(x) &= x^{\ln x} \cdot \frac{2 \ln x}{x} 
      \end{align*}
    \end{solution}

  \item
\begin{problem}[\stretch{1}]
      On what intervals is $f(x)$ increasing?  On what intervals is it
      decreasing?
    \end{problem}

    \begin{solution}
      \ideas{To answer this, let's make a sign chart for $f'$.}

      First, we'll find where $f'$ is 0 or undefined.  On $(0, \infty)$,
      $f'(x) = x^{\ln x} \cdot \frac{2 \ln x}{x}$ is always defined.
      Since $x^{\ln x}$ is always positive, $f'(x) = 0$ when $\ln x = 0$,
      which happens at $x = 1$.  Testing points gives us this sign chart:
      \begin{center}
        \begin{tikzpicture}
          \begin{axis}[axis y line=none, x axis line style={-}, xtick={0,1}, xmin=-1.5, xmax = 2, ymin=-2, ymax = 2, height=1in, width=3in, clip=false]
            \addplot+[domain=-1.5:2]{0};
            \draw (-1.5,1.5) node[right] {$f$};
            \draw (axis cs: -1.5, -1.5) node[right] {sign of $f'$};
            \draw (axis cs: 0.5, -1.5) node{$-$};
            \draw (axis cs: 1.5, -1.5) node{$+$};
            \draw[->] (0,1.5) -- (0.9,0.5);
            \draw[->] ( 0.9,0.5) -- (1.1,0.5);
            \draw[->] (1.1,0.5) -- (2,1.5);
          \end{axis}
        \end{tikzpicture}
      \end{center}
      So, \fbox{$f(x)$ is decreasing on $(0, 1)$ and increasing on $(1,
        \infty)$}.
    \end{solution}

  \item
    \begin{problem}[\stretch{0.2}]
      What is $\displaystyle \lim_{x \to \infty} x^{\ln x}$?

      \probonly{\emph{Hint: If $x$ is a large positive number, what is $\ln x$?  What about $x^{\ln x}$? What will happen as $x$ increases without bound?}}
    \end{problem}

    \begin{solution}
      As $x \to \infty$, $\ln x \to \infty$.  So, when $x$ is very large,
      $x^{\ln x}$ involves raising a big number to a big number, which
      gives a big result.  So, the limit is $\boxed{\infty}$.
    \end{solution}

  \item
    \begin{problem}[\stretch{0.2}]
      What is $\displaystyle \lim_{x \to 0^+} x^{\ln x}$?

      \probonly{\emph{Hint: What happens if $x$ is a positive number very close to $0$, like $x = 0.00001$?  What will happen as $x$ continues to approach $0$?}}
    \end{problem}

    \begin{solution}
      As $x \to 0^+$, $\ln x \to -\infty$, so when $x$ is close to 0,
      $x^{\ln x}$ involves raising a tiny positive number to a very
      negative power; the result is a big positive number.  So,
      $\displaystyle \lim_{x \to 0^+} x^{\ln x} = \boxed{\infty}$.
    \end{solution}

\pspace{\newpage}

  \item
    \begin{problem}[\stretch{1}]
      Sketch a rough graph of $f(x)$.  Your graph should accurately show
      where $f(x)$ is increasing and decreasing and reflect the limits you
      have found, but you need not show the concavity accurately.
    \end{problem}

    \begin{solution} % Jason Bland
      Here is the graph:
      \begin{center}
        \begin{tikzpicture}
          \begin{axis}[restrict y to domain=0:10, ymax=4, xtick={1}, ytick={1},height=6cm]
            \addplot+[domain=0:5, blue] {e^((ln(x))^2)};
          \end{axis}
        \end{tikzpicture}
      \end{center}
      The key features your graph should have are that it's decreasing on
      $(0, 1)$, increasing on $(1, \infty)$, and that both $\displaystyle
      \lim_{x \to 0^+} f(x)$ and $\displaystyle \lim_{x \to \infty} f(x)$
      are $\infty$. 
    \end{solution}

  \item
    \begin{problem}[\vfill]
      Does the graph of $f(x)$ have an global minimum on the interval $\left[
        \frac{1}{e^2}, e \right]$?  If so, where?  (Give both the $x$- and
      $y$-coordinate.)
    \end{problem}

    \begin{solution}
      From our sign chart in {{{{\#4}}(b)}}, we can see
      that the global minimum of $f$ on $\left[
        \frac{1}{e^2}, e \right]$ must happen at $x = 1$.  The
      corresponding $y$-coordinate is $f(1) = 1^{\ln 1} = 1^0 = 1$, so the
      global minimum point is $\boxed{(1, 1)}$. 
    \end{solution}

  \item
    \begin{problem}[\vfill]
      Does the graph of $f(x)$ have an global maximum on the interval $\left[
        \frac{1}{e^2}, e \right]$?  If so, where?  (Give both the $x$- and
      $y$-coordinate.)
    \end{problem}

    \begin{solution} % Jason Bland
      $f$ must have an global maximum on $\left[ \frac{1}{e^2}, e
      \right]$ by the Extreme Value Theorem: we're dealing with a
      continuous function on a closed interval.  From our sign chart in
      {{{{\#4}}(b)}}, we can see that the maximum must
      occur at an endpoint.  To decide which, we just evaluate $f$ at both
      enedpoints: $f(e) = e^1 = e$ while $f(\frac{1}{e^2}) =
      (\frac{1}{e^2})^{-2} = e^4$, and $e^4$ is larger. So the maximum is
      at the point $\boxed{\left( \frac{1}{e^2}, e^4 \right)}$.
    \end{solution}

  \end{enumerate}

\item 
\begin{grading}
3 points. 1 point for identifying B as the best strategy, 1 point for identifying C as possible, and 1 point for identifying A as unworkable. 
\end{grading}

 \textbf{Reflect Back.}

  \begin{problem}[\stretch{1}]
    To differentiate $y =10 (x^2 + 7)^{\pi} +
    \displaystyle \frac{ \sqrt 2(2x^2 + 3)^{5}}{3} + 
    \frac{x^8(3x^2+2)^{4x}}{x^2+1}$, which of the following would be the most efficient strategy?
    Which would work, but would take some energy?  Which is totally unworkable?
    \begin{enumerate}[label=\protect\maybecircle{\Alph*.}]
    \item Take the natural log of both sides and use logarithmic
      differentiation.
    \item Differentiate the
      first term using the Chain Rule.  For the second term,
      first pull out $\frac{\sqrt 2}{3}$ as a constant, and then
      differentiate $(2x^2 + 3)^{5}$ using the Chain Rule.
      Then differentiate the third term using logarithmic differentiation.
    \item Differentiate each of
      the three terms using logarithmic differentiation.
    \end{enumerate}
  \end{problem}

  \begin{solution}
    (A) would not work; we can't simplify the natural log of a sum.  (B) and (C)
    both work, but (B) is a bit more efficient.
  \end{solution}

\pspace{\newpage}

\item 
\begin{problem}
    We are going to start our unit on trigonometry next class. To prepare, let's review some of the geometric properties of circles.

    In the United States, small pizzas are usually 10 inches in diameter, and medium pizzas are usually 12 inches in diameter.
\end{problem}

\begin{enumerate}
    \item 
    \begin{problem}
    Your friend buys pizzas just to eat the crusts. How much crust (in inches) is on a small pizza? A medium pizza?
    \end{problem}
    \pspace{\vfill}

    \begin{solution}
        The length of crust is the circumference of a circle. We know that $C=2\pi r$, where $r$ is the radius of the circle.
        \begin{itemize}
            \item A small pizza (radius of 5 inches) has $10\pi$ inches of crust.
            \item A medium pizza (radius of 6 inches) has $12\pi$ inches of crust.
        \end{itemize}
    \end{solution}

    \item 
    \begin{problem}
    Suppose we slice the pizzas into six equal slices. How many slices of a medium pizza would give the same amount of crust as an entire small pizza?
    \end{problem}
    \pspace{\vfill}

    \begin{solution}
        If we slice a medium pizza into six equal slices, then each slice will have $2\pi$ inches of crust. 5 of these slices would give the same amount of crust as an entire small pizza.
    \end{solution}

    \item 
    \begin{problem}
        Your local pizzeria offers a 6-inch-diameter pizza at half the price of a medium pizza. Do you think this is a fair deal? Would your friend think this is a fair deal? Explain your reasoning.
    \end{problem}
    \pspace{\vfill}

    \begin{solution}
        A 6-inch diameter pizza has a circumference of $6\pi$ inches and an area of $9\pi$ square inches. A medium pizza (12-inch diameter) has a circumference of $12\pi$ inches and an area of $36\pi$ square inches. If we care about the total amount of pizza we're getting and not just the crust, then this isn't a good deal, because we are getting a fourth as much pizza for half the cost. If we only care about the crust, like our friend, then we might say this is a fair deal, because we are getting half as much crust for half the cost.
    \end{solution}
\end{enumerate}

\end{enumerate}

\psreflection{

}

\end{document}